\documentclass[11pt,a4paper]{article}
\usepackage[T1]{fontenc}
\usepackage[utf8]{inputenc}
\usepackage[french]{babel}
\usepackage{lmodern}
\usepackage{amsmath,amssymb}
\usepackage[margin=2.35cm]{geometry}
\usepackage{booktabs,longtable,tabularx}
\usepackage{enumitem}
\usepackage{xcolor}
\usepackage{listings}
\usepackage{microtype}
\usepackage{hyperref}

\hypersetup{colorlinks=true,linkcolor=blue!45!black,urlcolor=blue!45!black}
\setlist{nosep}
\newcommand{\code}[1]{\texttt{#1}}
\newcommand{\NW}{\mathrm{NW}}
\newcommand{\E}{\mathbb{E}}
\newcommand{\R}{\mathbb{R}}
\definecolor{codebg}{RGB}{246,247,249}
\lstset{basicstyle=\ttfamily\small,backgroundcolor=\color{codebg},
  frame=single,breaklines=true,columns=fullflexible,showstringspaces=false}

\title{Mécanique du modèle de société de crédit M4\\
\large et spécification de sa réduction minimale M4B}
\author{Audit et réimplémentation du dossier \code{m4\_credit\_soc\_fable}}
\date{17 juillet 2026}

\begin{document}
\maketitle

\begin{abstract}
Ce document décrit la mécanique effectivement exécutée par le modèle M4,
et non les variantes historiques qui restent accessibles pour les ablations.
Chaque entité ne possède qu'un stock réel, son capital productif $K$. Les
créances et dettes proviennent exclusivement d'un réseau de prêts nominaux
perpétuels. À chaque pas, des naissances et des chocs alimentent le système,
la production accroît $K$, les intérêts redistribuent du capital, le marché
crée de nouvelles expositions, puis les défauts détruisent les créances et
peuvent déclencher une cascade. M4B fixe les institutions de la configuration
finale validée, supprime les branches expérimentales et conserve les mêmes
tirages aléatoires et le même ordre de calcul. Des tests de parité comparent
les deux moteurs à chaque pas et obtiennent des trajectoires identiques.
\end{abstract}

\tableofcontents

\section{Ce que le système représente}

M4 est un modèle multi-agents de société de crédit. Il n'y a ni banque
centrale, ni monnaie séparée, ni bilan bancaire agrégé. Les entités naissent,
produisent, prêtent une partie de leur capital à d'autres entités et paient
des intérêts sur leurs dettes. Leur réseau de contrats produit une fragilité
endogène : la disparition d'une emprunteuse efface l'actif de ses prêteuses.

La configuration principale est volontairement asymétrique :

\begin{itemize}
\item $K$ est réel, productif, soumis aux chocs et à la dépréciation ;
\item le principal d'un prêt est nominal, perpétuel et ne se déprécie pas ;
\item seuls les intérêts sont servis ; le principal n'est jamais amorti ;
\item une entité sans dette contractuelle ne peut pas mourir ;
\item aucune récupération ne protège une prêteuse lorsque sa débitrice meurt.
\end{itemize}

Cette asymétrie est le cœur du mécanisme. Le marché reconstruit sans cesse un
stock nominal de créances adossé à des capitaux réels fluctuants. Les pertes
sur créances sont donc capables de faire passer une prêteuse de solvable à
insolvable sans qu'un choc exogène supplémentaire soit nécessaire.

\section{État du modèle}

\subsection{Entités}

Une entité $i$ est identifiée par un entier attribué à sa naissance et jamais
réutilisé. Son seul état économique réel est
\[
K_i(t)\geq 0.
\]
Le moteur mémorise aussi sa date de naissance, son statut vivant ou mort et,
pour le pas courant, sa production, ses intérêts reçus, ses intérêts payés et
un indicateur de défaut de service. Ces dernières quantités sont des mesures,
pas des stocks reportés au pas suivant.

L'ensemble des identifiants vivants est maintenu explicitement et trié avant
chaque parcours. Ce détail garantit que le même vecteur de chocs aléatoires
est toujours affecté aux mêmes identifiants pour une graine donnée.

\subsection{Contrats de prêt}

Un prêt est un quadruplet
\[
(\ell,b,q,r),
\]
où $\ell$ est la prêteuse, $b$ l'emprunteuse, $q>0$ le principal et $r>0$
le taux d'intérêt par pas. Le contrat est perpétuel. Pour chaque entité, le
carnet maintient les agrégats
\[
C_i=\sum_{(i,b,q,r)}q,
\qquad
D_i=\sum_{(\ell,i,q,r)}q,
\qquad
S_i=\sum_{(\ell,i,q,r)}rq.
\]
$C_i$ est le total de ses créances, $D_i$ sa dette nominale et $S_i$ le
service d'intérêts dû au pas courant.

Il existe au plus un contrat orienté par paire $(\ell,b)$. Si la même paire
échange de nouveau, les principaux sont additionnés et le taux devient
\[
r_{\mathrm{new}}
=\frac{q_0r_0+qr}{q_0+q}.
\]
Cette fusion préserve exactement le principal total et le flux d'intérêts
$q_0r_0+qr$. La paire inverse $(b,\ell)$ reste un contrat différent.

La valeur nette d'une entité est dérivée du capital et du carnet :
\[
\boxed{\NW_i=K_i+C_i-D_i.}
\]
Il n'existe aucun terme de dette initiale $d_0$. Une entité qui n'a jamais
emprunté vérifie $D_i=0$ et ne peut donc pas devenir insolvable tant que
$K_i\geq0$.

\section{Chronologie exacte d'un pas}

L'ordre des opérations fait partie du modèle. Une permutation des phases
changerait la date du premier paiement d'un prêt, l'ensemble admis sur le
marché ou le déclenchement des faillites.

\begin{center}
\begin{tabularx}{\textwidth}{@{}clX@{}}
\toprule
Phase & Opération & Effet principal \\
\midrule
1 & Naissances & injection de $K_0$ et nouveaux identifiants \\
2 & Choc & variation multiplicative du capital réel \\
3 & Production & création réelle, entièrement retenue dans $K$ \\
4 & Intérêts & transfert de $K$ ; défaut si liquidité insuffisante \\
5 & Dépréciation & destruction d'une fraction du capital réel \\
6 & Marché & transfert de $K$ et création de contrats nominaux \\
7 & Faillites & annulation de contrats, destruction et propagation \\
8 & Mesure & séries, avalanches et instantanés, sans tirage aléatoire \\
\bottomrule
\end{tabularx}
\end{center}

\subsection{Phase 1 : naissances}

Le nombre de naissances est
\[
B_t\sim\operatorname{Poisson}(\lambda).
\]
Chaque nouvelle entité reçoit $K_0$ et aucun contrat. L'injection réelle du
pas est donc $I_t=B_tK_0$. Les nouvelles entités participent immédiatement
aux phases suivantes : elles subissent le choc, produisent, puis peuvent
entrer sur le marché du même pas.

Si la population finit un pas à zéro, le pilote prend l'état
\code{extinction} et s'arrête. L'extinction est donc absorbante dans une
exécution normale, même si une loi de Poisson aurait permis des naissances
ultérieures.

\subsection{Phase 2 : choc multiplicatif}

Pour chaque vivante,
\[
K_i\leftarrow K_i\exp(\eta_i),
\qquad
\eta_i\overset{\mathrm{iid}}{\sim}
\mathcal N\!\left(-\frac{\sigma^2}{2},\sigma^2\right).
\]
Le décalage $-\sigma^2/2$ donne
\[
\E[\exp(\eta_i)]=1.
\]
Le choc est ainsi sans dérive en espérance, bien que sa médiane soit
inférieure à un. La grandeur enregistrée
\code{shock\_gain} est la différence entre le capital total juste après et
juste avant cette phase ; elle peut être positive ou négative dans un run.

\subsection{Phase 3 : production}

La production est concave :
\[
P_i=\alpha\sqrt{K_i},\qquad \alpha=1.
\]
Elle est entièrement retenue :
\[
K_i\leftarrow K_i+P_i.
\]
Il n'y a plus de variable de liquidité, de taux d'épargne ou de consommation
séparée. Le champ \code{prod} conserve $P_i$ jusqu'au pas suivant afin de
calculer le revenu individuel dans les instantanés.

\subsection{Phase 4 : service des intérêts}

Pour une emprunteuse $b$, le montant dû est $S_b$. Si $K_b\geq S_b$, elle
paie tout. Sinon, elle verse tout son capital au prorata et fait défaut :
\[
\phi_b=\min\left(1,\frac{K_b}{S_b}\right),
\qquad
p_{\ell b}=\phi_b r_{\ell b}q_{\ell b}.
\]
Après paiement,
\[
K_b\leftarrow K_b-\sum_\ell p_{\ell b},
\qquad
K_\ell\leftarrow K_\ell+p_{\ell b}.
\]
Si $\phi_b<1$, $K_b$ devient nul et son indicateur de défaut est activé. Le
principal et le service futur du contrat restent inchangés : il n'y a ni
arriéré, ni renégociation. Le défaut entraînera la mort en phase 7, même si
la valeur nette comptable de l'entité reste positive.

Une subtilité est conservée pour la fidélité exacte : les emprunteuses sont
servies selon l'ordre historique d'apparition des clés du carnet. Une entité
qui reçoit des intérêts comme prêteuse avant son propre tour peut employer
ce capital pour payer. Ce service est proportionnel entre les créancières
d'une même emprunteuse, mais il n'est pas simultané à l'échelle du réseau.

Un prêt créé pendant la phase 6 ne paiera ses premiers intérêts qu'au pas
suivant. Cette temporalité évite de faire payer un contrat avant sa création.

\subsection{Phase 5 : dépréciation}

Le capital réel devient
\[
K_i\leftarrow(1-\delta)K_i.
\]
Une valeur résiduelle inférieure à $10^{-12}$ est arrondie à zéro. Les
principaux $q$ et les taux $r$ sont nominaux : ils restent strictement
inchangés. La dépréciation fragilise donc mécaniquement un bilan endetté en
réduisant son actif réel sans réduire son passif nominal.

\subsection{Phase 6 : marché local du crédit}

Le pool est la liste des entités vivantes qui n'ont pas fait défaut pendant
le service du pas. Il est figé pour toute la phase. Le moteur réalise un
round par membre du pool. À chaque round :

\begin{enumerate}
\item il tire uniformément jusqu'à $k$ identifiants distincts ;
\item la plus riche en $K$ devient prêteuse $\ell$, la plus pauvre devient
  emprunteuse $b$ ;
\item si leurs capitaux diffèrent, il calcule un taux et une cible commune ;
\item il transfère le volume possible et ajoute le contrat correspondant.
\end{enumerate}

Le rendement marginal imputé à une entité de capital $K_x$ est
\[
r_x=\frac{\alpha}{2\sqrt{K_x}}.
\]
Le taux bilatéral est la moyenne géométrique
\[
r=\sqrt{r_\ell r_b}.
\]
L'objectif de revenu définit la cible
\[
K^*(r)=\left(\frac{\alpha}{2r}\right)^2
=\sqrt{K_\ell K_b}.
\]
Le principal transféré vaut
\[
q=\min\!\left(K_\ell-K^*,\;\max(0,K^*-K_b)\right).
\]
Si $q\geq10^{-9}$,
\[
K_\ell\leftarrow K_\ell-q,
\qquad
K_b\leftarrow K_b+q,
\]
et le contrat $(\ell,b,q,r)$ est inscrit. Cette opération conserve le
capital réel total. Elle conserve aussi la valeur nette de chacune des deux
parties au moment de l'échange : la prêteuse remplace $q$ de capital par une
créance de $q$, et l'emprunteuse reçoit $q$ de capital contre une dette de
$q$.

Pour les grandes populations, le tirage sans remise utilise un rejet de
doublons plutôt qu'une permutation complète. Ce choix est une optimisation,
mais sa branche et son ordre d'appels au générateur aléatoire sont conservés
dans M4B afin d'obtenir la même trajectoire que M4.

\subsection{Phase 7 : faillites et cascades}

Une première file, triée par identifiant, rassemble toutes les vivantes qui
vérifient au moins une condition :
\[
\text{défaut de service}
\qquad\text{ou}\qquad
\NW_i<-10^{-9}.
\]
La cause d'une racine est notée \code{liquidity}, \code{insolvency} ou
\code{both}. Le traitement d'une faillite $i$ est irréversible :

\begin{enumerate}
\item chaque contrat où $i$ est emprunteuse est supprimé ; la prêteuse perd
  son principal et une arête de perte $i\rightarrow\ell$ de poids $q$ est
  enregistrée ;
\item chaque contrat où $i$ est prêteuse est supprimé ; ses emprunteuses sont
  libérées de ces dettes ;
\item le capital $K_i$ restant est détruit ; $i$ est retirée de la population
  vivante et son capital est mis à zéro.
\end{enumerate}

Après le traitement de toute la file, le moteur recalcule la valeur nette des
survivantes. Celles qui sont devenues insolvables par perte d'une créance
forment la génération suivante. Le processus continue jusqu'à une file vide.
La profondeur enregistrée est le numéro maximal de génération, et non la
longueur d'un chemin particulier dans le graphe.

Les membres d'une même file sont déterminés avant le début de l'itération.
Une entité déjà placée dans cette file sera donc tuée même si l'annulation
d'une de ses propres dettes par une faillite précédente améliorerait ensuite
son bilan. Cette convention de mise à jour par générations est identique dans
M4 et M4B.

\subsection{Phase 8 : mesures}

Une ligne de série est construite après la résolution complète. Elle contient
la démographie, le capital et la valeur nette totaux, la production, le
nombre et le volume de prêts, les paiements, les défauts, les causes racines,
la profondeur, les pertes de créances, le capital détruit et les trois flux
du bilan réel. Cette phase ne tire aucun nombre aléatoire et ne modifie pas
la trajectoire économique.

\section{Construction d'une avalanche}

Toutes les morts d'un même pas ne constituent pas automatiquement une seule
avalanche. Le moteur construit le graphe des arêtes de perte créées pendant
la résolution, puis le restreint aux entités mortes lors de cette résolution.
Une avalanche est une composante faiblement connexe de ce graphe.

Ainsi, deux racines sans exposition commune restent deux avalanches de taille
un. Deux racines déjà insolvables mais liées par un contrat appartiennent en
revanche à la même composante ; le champ \code{n\_roots} permet de distinguer
ce cas d'une longue propagation. Les mesures fondamentales sont :

\begin{itemize}
\item \code{size} : nombre de mortes dans la composante ;
\item \code{n\_roots} : nombre de membres présents dans la première file ;
\item \code{depth} : génération maximale atteinte ;
\item \code{generation} : génération de chaque membre ;
\item \code{volume\_j} : somme des principaux de créances détruits par les
faillites de la composante, y compris lorsque la prêteuse touchée survit ;
\item le rapport de branchement empirique
\[
b=1-\frac{\sum_a R_a}{\sum_a S_a},
\]
où $R_a$ et $S_a$ sont le nombre de racines et la taille de l'avalanche $a$.
\end{itemize}

La mesure retenue est strictement intra-pas. Relier automatiquement des
pertes du pas $t$ aux morts du pas $t+1$ peut fabriquer une composante géante
par percolation temporelle ; cette variante exploratoire n'est donc pas
intégrée au moteur minimal.

\section{Pourquoi des cascades apparaissent}

Le mécanisme est une alternance accumulation--relaxation.

Pendant les périodes calmes, chaque pas ajoute environ un round de marché par
entité. Le carnet fusionne les échanges répétés, mais le stock d'expositions
nominales est continuellement reconstruit. Les chocs et la dépréciation font
évoluer le capital réel sous ce réseau nominal. Une entité peut alors faire
défaut par manque de capital disponible ou devenir insolvable.

Sa disparition retire brutalement $q$ de l'actif comptable de chacune de ses
prêteuses. Il n'existe dans la règle finale ni récupération du capital
résiduel, ni transfert des créances détenues par la morte. Si la perte rend
une prêteuse insolvable, celle-ci disparaît à la génération suivante et
transmet à son tour ses propres pertes. La cascade élague ainsi simultanément
la population et le réseau, avant que les naissances et le marché ne le
reconstruisent.

Le volume de marché est donc une propriété institutionnelle déterminante. La
configuration finale fixe un round par tête : le paramètre expérimental qui
divisait ou multipliait ce volume n'est pas conservé dans M4B.

\section{Invariants et contrôles comptables}

\subsection{Capital non négatif}

Le choc conserve le signe, la production est positive, le paiement est borné
par le capital disponible, le marché ne prélève pas plus que l'excédent à la
cible et la mort fixe $K$ à zéro. Par construction,
\[
K_i\geq0.
\]

\subsection{Égalité des créances et des dettes}

Chaque principal apparaît une fois comme créance et une fois comme dette :
\[
\sum_i C_i=\sum_i D_i.
\]
Après chaque cascade, aucun contrat ne doit référencer une entité morte. Il
en découle, à l'arrondi près,
\[
\sum_{i\ \mathrm{vivante}}\NW_i
=\sum_{i\ \mathrm{vivante}}K_i.
\]

\subsection{Bilan réel d'un pas}

Les prêts et les intérêts sont des transferts internes. Les seuls flux qui
changent le capital total sont l'injection des naissances, le choc réalisé,
la production, la dépréciation et la destruction lors des faillites :
\[
\boxed{
K_{\mathrm{tot}}(t)-K_{\mathrm{tot}}(t-1)
=I_t+G_t+P_t-\Delta_t-X_t.}
\]
Les colonnes correspondantes de \code{series.csv} permettent de tester cette
identité après coup.

\subsection{Reproductibilité}

Chaque simulation possède un unique générateur
\code{numpy.random.default\_rng(seed)}. L'ordre est : Poisson des naissances,
vecteur normal des chocs, puis échantillons du marché. Les faillites et les
mesures ne consomment aucun tirage. Les instantanés accèdent aux agrégats avec
des lectures non mutantes, afin de ne pas insérer de nouvelles clés dans les
tables ordonnées du carnet. Un run enregistré et un run non enregistré ont
donc la même série bit à bit.

\section{Ce que M4B retire et ce qu'il conserve}

M4 contient des variantes utiles pour établir la causalité scientifique, mais
elles ne décrivent pas le modèle principal. Les laisser dans une version
pédagogique ferait croire que toutes appartiennent à la mécanique finale.

\begin{longtable}{p{0.29\textwidth}p{0.29\textwidth}p{0.34\textwidth}}
\toprule
Élément & Décision M4B & Motif \\
\midrule
\endhead
État réel $K$ & conservé & unique stock économique du modèle final \\
Prêts perpétuels et fusion par paire & conservés et fixés & structure
constitutive du réseau \\
Choc i.i.d. & conservé et fixé & configuration principale ; pas de
corrélation imposée \\
Objectif de revenu & conservé et fixé & donne $K^*=\sqrt{K_\ell K_b}$ \\
Un round par tête & conservé et fixé & institution finale validée \\
Annulation + destruction & conservées et fixées & mécanisme de propagation \\
Chocs macro et sectoriels & retirés & ablations sur la synchronisation \\
Objectif de richesse & retiré & contrôle expérimental à faible levier \\
Naissance par emprunt & retirée & variante non active \\
Cohorte initiale & retirée & ablation démographique \\
Désactivation du crédit & retirée & réduction de contrôle, pas modèle principal \\
Transfert des créances / récupération prorata & retirés & anciennes règles
qui amortissent les cascades \\
Nombre de rounds configurable & retiré & levier de campagne fixé à sa valeur
finale \\
Logs enrichis, watchers, figures, ajustements dPlN & séparés & post-traitement,
sans rôle dynamique \\
\bottomrule
\end{longtable}

Les constantes numériques de tolérance et $\alpha=1$ sont définies une seule
fois dans le module, et non présentées comme des paramètres de recherche. La
configuration publique ne contient que :

\begin{center}
\begin{tabular}{@{}lll@{}}
\toprule
Nom & Défaut & Rôle \\
\midrule
\code{lam} & 10 & intensité des naissances \\
\code{delta} & 0,05 & dépréciation du capital \\
\code{sigma} & 0,25 & volatilité du choc \\
\code{K0} & 25 & capital de naissance \\
\code{k} & 3 & taille de l'échantillon de marché \\
\code{seed} & 0 & graine du générateur \\
\code{T} & 2000 & horizon maximal \\
\code{pop\_max} & 30000 & garde-fou d'arrêt \\
\bottomrule
\end{tabular}
\end{center}

\section{Données enregistrées pour l'analyse statistique}

Le moteur et l'écriture sont séparés : \code{model.py} ne connaît aucun
chemin de fichier, tandis que \code{io.py} reçoit des instantanés en lecture
seule. Un dossier de run contient les données primaires suivantes.

\begin{longtable}{p{0.31\textwidth}p{0.61\textwidth}}
\toprule
Fichier & Contenu et usages \\
\midrule
\endhead
\code{config.json} & version, paramètres, $\alpha$ et règles institutionnelles
fixes ; reproduction exacte du run \\
\code{series.csv} & une ligne par pas : démographie, capital, crédit, défauts,
avalanches et bilan réel \\
\code{individual\_series.csv.gz} & une ligne par entité vivante et par pas,
plus une ligne terminale pour chaque morte ; trajectoires individuelles
compressées \\
\code{avalanches.csv} & identifiant, pas, taille, profondeur, nombre et causes
des racines, volume causal des pertes de créances en joules \\
\code{avalanche\_members.csv} & appartenance de chaque morte, statut racine et
génération ; audit des agrégats causaux \\
\code{entities.csv} & naissance, mort éventuelle, cause et survie finale de
chaque identifiant \\
\code{deaths.csv} & âge et bilan individuel juste avant chaque mort \\
\code{final\_loans.csv} & prêteuse, emprunteuse, principal et taux du carnet
final \\
\code{entities\_t*.npz} & vecteurs individuels $K,C,D,\NW$,
production, revenu, âge et degrés \\
\code{network\_t*.npz} & arêtes $(\ell,b,q,r)$ aux mêmes dates \\
\code{summary.json} & statut final, comptes globaux, rapport de branchement et
résultat du contrôle de cohérence \\
\bottomrule
\end{longtable}

Ces tables suffisent à étudier la stationnarité de la population, la
distribution des tailles et profondeurs d'avalanches, le rapport de
branchement, les durées de vie, les distributions de $K$, $\NW$ et du revenu,
le renouvellement du haut de la distribution et la topologie du réseau. Les
figures et ajustements statistiques restent des produits dérivés : ils
peuvent être régénérés sans relancer le modèle.

Le volume d'une avalanche est affecté par la source de chaque perte : pour
une composante $a$, $V_a$ additionne tous les principaux dont l'emprunteuse
faillie appartient à $a$. Cette convention conserve donc la perte infligée à
une prêteuse survivante et permet les CCDF et ajustements dPlN sans pooling
ambigu de plusieurs avalanches survenues au même pas.

La table longitudinale individuelle contient précisément : $t$, identifiant,
date de naissance, âge, statut de vie, cause et génération de décès, $K$,
créances, dettes, $\NW$, production, intérêts reçus et payés, revenu brut,
revenu net, défaut de service et degrés entrant/sortant du réseau. Pour une
survivante, le bilan est celui de la fin du pas, après la cascade. Pour une
morte, la ligne terminale conserve son bilan juste avant la suppression de
ses contrats et de son capital. Toutes les entités possèdent donc une
trajectoire jusqu'à leur dernier état économiquement interprétable.

Cette table est écrite en flux dans un CSV compressé : sa taille ne pèse pas
sur la mémoire de la simulation. Elle est exhaustive par défaut. La fréquence
des survivantes peut être réduite pour les très grandes populations ; l'état
terminal de chaque morte et le dernier état des survivantes restent alors
conservés. Cette option ne change que la mesure et ne consomme aucun tirage
aléatoire.

Les instantanés sont espacés de 50 pas par défaut. Une valeur nulle désactive
les instantanés intermédiaires, mais l'état final est toujours enregistré.
Le format NumPy compressé conserve les types et évite la duplication d'une
colonne de temps sur chaque ligne individuelle.

\section{Validation de la réduction}

La fidélité n'est pas évaluée seulement sur des moyennes. Le test de parité
lance M4 et M4B avec la même configuration principale et, après chaque pas,
compare exactement :

\begin{itemize}
\item le statut et toutes les colonnes communes de la série ;
\item le capital, la vie, la naissance et le défaut de chaque entité ;
\item tous les contrats, créances, dettes et services dus ;
\item les tailles, profondeurs, racines et causes des avalanches.
\end{itemize}

Deux régimes, avec graines, volatilités, intensités de naissance et tailles
d'échantillon différentes, couvrent 210 pas dans le test automatisé. Les
trajectoires sont identiques champ par champ. Cinq contrôles supplémentaires
vérifient le bilan réel, l'absence de contrats orphelins, le paiement partiel,
la neutralité des mesures, la fréquence des trajectoires individuelles et la
relecture des fichiers de sortie.

La suite de tests du M4 de référence a également été exécutée avant la
réduction : 69 tests sur 69 passent. Un condensat n'est donc pas construit à
partir d'une référence déjà défaillante.

\section{Utilisation}

Depuis le dossier \code{m4b\_credit\_soc\_mini} :

\begin{lstlisting}[language=bash]
python3 run.py --steps 2000 --seed 0 --output results/baseline
\end{lstlisting}

Les principaux paramètres peuvent être changés explicitement :

\begin{lstlisting}[language=bash]
python3 run.py --steps 4000 --lambda 30 --sigma 0.25 \
  --market-sample 3 --snapshot-every 100 --seed 2 \
  --output results/lambda30_seed2
\end{lstlisting}

Le test autonome se lance avec :

\begin{lstlisting}[language=bash]
python3 tests/test_mini.py
\end{lstlisting}

La dépendance d'exécution est limitée à NumPy. SciPy, les bibliothèques de
tracé et les routines d'ajustement de lois ne sont pas nécessaires pour
simuler ou enregistrer un run.

\section{Traçabilité du code audité}

La lecture a porté sur le rapport M4 et sur tous les modules dynamiques. Les
points d'ancrage principaux sont :

\begin{itemize}
\item configuration : \code{m4\_credit\_soc\_fable/src/m4/config.py} ;
\item entités et contrats : \code{entities.py}, \code{contracts.py} ;
\item chocs, production et dépréciation : \code{production.py} ;
\item marché : \code{market.py} ;
\item faillites et avalanches : \code{bankruptcy.py} ;
\item ordre d'exécution et mesures : \code{simulation.py} ;
\item schémas historiques de sortie : \code{metrics.py} ;
\item justification scientifique :
  \code{reports/01\_soc\_final/main.tex} et \code{NOTES.md} ;
\item tests de la référence : \code{tests/}.
\end{itemize}

Le dossier de référence a été traité comme immuable. M4B est une copie
conceptuelle indépendante : il n'importe aucun module de M4 pendant une
simulation normale. Seul le test de parité importe les deux moteurs afin de
les comparer.

\section*{Conclusion}
\addcontentsline{toc}{section}{Conclusion}

La mécanique finale se résume à un stock réel productif, un réseau nominal
perpétuel et une règle de perte sans amortisseur. Le marché accumule des
expositions ; les chocs, la dépréciation et le service des intérêts déplacent
les bilans ; une faillite efface des créances et peut rendre insolvables les
prêteuses exposées. M4B conserve exactement cette boucle et son ordre de
calcul, tout en retirant les axes d'ablation et l'infrastructure de campagne.
La séparation entre moteur, données primaires et analyse rend désormais le
système plus court à lire, plus difficile à paramétrer accidentellement hors
du modèle validé et directement exploitable pour de futures statistiques.

\end{document}
